\begin{theorem}[Pointwise stream metric regular-Cauchy pruning route]
\label{thm:pointwise-stream-metric-regular-cauchy-pruning-route}
Every completion-facing $\PointwiseStreamMetricUp$ read that prunes redundant pointwise observations must first expose the $\StreamNameUp$ pointwise windows, $\RegSeqRatUp$ readbacks, $\DyadicRatCoreUp$ tolerance ledger, and terminal $\RealUp$ seal already carried by \autoref{def:pointwise-stream-metric-carrier}. The pruning handoff then factors through the $\RegularCauchyPruningUp$ certificate of \autoref{ch:concrete-instances-regular-cauchy-pruning-namecert} before any downstream metric or completion consumer reads the retained stream. The route exports no host stream equality, hidden infinite selector, quotient Real equality, or ambient completeness principle.
\end{theorem}
\begin{proof}
Project the pointwise stream metric carrier. Its analytic rows expose the stream source, finite pointwise windows, dyadic tolerance ledger, regular rational readback, Real seal, and metric row before any completion-facing consumer has a pointwise distance packet to inspect.

Apply the regular-Cauchy pruning certificate of \autoref{ch:concrete-instances-regular-cauchy-pruning-namecert}. That sibling consumes exactly the StreamName windows, retained schedule, dyadic ledger, RegSeqRat readback, and inherited Real seal needed to delete redundant observation requests as finite NameCert data.

Finally use the existing L10 route of \autoref{thm:pointwise-stream-metric-l10-route}. The retained stream remains inside the displayed pointwise windows and Real seal boundary before it reaches metric or completion consumers, so the horizontal edge is a concrete stream/pruning dependency rather than a kernel-only pointwise metric export.
\end{proof}
